%! Author = boldizsar.poor
%! Date = 26/10/2026

% Preamble
\documentclass[11pt]{article}

% Packages
\usepackage{amsmath}
\usepackage{stmaryrd}
\usepackage{csquotes}
\usepackage{bookmark}
\usepackage{xcolor}

\usepackage{amsthm,thmtools}
\theoremstyle{plain}
\newtheorem{theorem}{Theorem}[section]
\newtheorem{lemma}[theorem]{Lemma}
\newtheorem{proposition}[theorem]{Proposition}\theoremstyle{definition}
\newtheorem{definition}[theorem]{Definition}
\newtheorem{example}[theorem]{Example}
\newtheorem{remark}[theorem]{Remark}

\usepackage[noabbrev,nameinlink]{cleveref}
\usepackage{todonotes}

\usepackage{macros/tikzfig}
\input{macros/floquet.tikzstyles}
\input{macros/floquet.tikzdefs}

\usepackage[backend=biber]{biblatex}
\addbibresource{macros/AdaptiveSteane.bib}

% Document
\begin{document}

\title{CSS Fault Tolerance}

\maketitle

\tableofcontents

\section{Gottesman's Definition of Fault-tolerance\texorpdfstring{~\cite{gottesmanIntroductionQuantumError2009,gottesmanSurvivingQuantumComputer2026}}{}}

\begin{definition}
    The \emph{$r$-filter} for the QECC $Q$ is the projector on the subspace spanned by states of the form $E \ket{\psi}$, where $\ket{\psi} \in Q$ and $E$ is an error of weight at most $r$.
    Pictorially, I will draw an $r$-filter as:
    \[
        \tikzfig{figures/Gottesman/r-filter.tikz}
    \]
    The doubled horizontal lines in this picture represent a block of the QECC $Q$ rather than single physical qubits.
\end{definition}
\begin{definition}
    The \emph{ideal decoder} for the QECC $Q$ is the quantum channel that takes a state (possibly with errors) encoded in $Q$, corrects the errors, and then decodes the logical state, discarding all ancilla qubits.
    I will draw an ideal decoder as follows:
    \[
        \tikzfig{figures/Gottesman/ideal-decoder.tikz}
    \]
    The doubled horizontal line on the left of the picture represents a block of $Q$.
    The single horizontal line on the right of the picture is the logical qubit after decoding.
\end{definition}

\subsection{Fault-tolerant Gate Properties}

\begin{definition}
    The picture to represent a noisy gate gadget with exactly $s$ faults is
    \[
        \tikzfig{figures/Gottesman/noisy-1-gate.tikz}
        \qquad
        \text{or}
        \qquad
        \tikzfig{figures/Gottesman/noisy-2-gate.tikz}
    \]
    for one- and two-qubit gates, respectively.
    Each doubled horizontal line represents a single block of the QECC\@.

    The picture to represent an ideal gate performed on unencoded qubits without noise is
    \[
        \tikzfig{figures/Gottesman/ideal-1-gate.tikz}
        \qquad
        \text{or}
        \qquad
        \tikzfig{figures/Gottesman/ideal-2-gate.tikz}
    \]
    for one- and two-qubit gates, respectively. Each single horizontal line represents a single unencoded qubit.
\end{definition}

\begin{definition}[FT Gate Error Propagation Property]
    Suppose we have a gate gadget associated with a QECC with distance $2t + 1$ (i.e., it corrects $t$-errors).
    If it is a single-qubit gate gadget, it satisfies the \emph{FT Gate Error Propagation Property}, abbreviated \emph{GPP}, if, whenever $r + s \leq t$,
    \[
        \tikzfig{figures/Gottesman/GPP1.tikz}
    \]
    A two-qubit gate gadget satisfies the GPP if, whenever $r_1 + r_2 + s \leq t$,
    \[
        \tikzfig{figures/Gottesman/GPP2.tikz}
    \]
    For both equations, the faulty locations on both sides of the equation have the same fault path and error maps substituted in.
\end{definition}
\begin{definition}[FT Gate Correctness Property]
    Suppose we have a gate gadget associated with a QECC with distance $2t + 1$.
    If it is a single-qubit gate gadget, it satisfies the FT Gate Correctness Property, abbreviated GCP, if, whenever $r + s \leq t$,
    \[
        \tikzfig{figures/Gottesman/GCP1.tikz}
    \]
    A two-qubit gate gadget satisfies the GCP if, whenever $r_1 + r_2 + s \leq t$,
    \[
        \tikzfig{figures/Gottesman/GCP2.tikz}
    \]
    These equations must hold for any choice of CP maps substituted in for the faulty locations on the left-hand side.
\end{definition}

\subsection{Fault-tolerant State Preparation}

\begin{definition}
    The picture for a single-qubit state preparation gadget with exactly $s$ faults is
    \[
        \tikzfig{figures/Gottesman/noisy-state-prep.tikz}
    \]
    The picture for a single-qubit measurement gadget with exactly $s$ faults is
    \[
        \tikzfig{figures/Gottesman/noisy-measurement.tikz}
    \]
    In both cases, the doubled horizontal lines represent encoded blocks of the QECC\@.
    The measurement gadget outputs classical information, generally a single bit.

    The pictures for ideal preparation of a single-qubit state and ideal single-qubit measurement are
    \[
        \tikzfig{figures/Gottesman/ideal-state-prep.tikz}
        \qquad
        \text{and}
        \qquad
        \tikzfig{figures/Gottesman/ideal-measurement.tikz}
    \]
    The horizontal lines represent single unencoded qubits, and the ideal measurement outputs a classical bit.
\end{definition}

\begin{definition}[FT Preparation Error Propagation Property]
    Suppose we have a preparation gadget associated with a QECC with distance $2t + 1$.
    The gadget satisfies the \emph{FT Preparation Error Propagation Property}, abbreviated \emph{PPP}, if, whenever $s \leq t$,
    \[
        \tikzfig{figures/Gottesman/PPP.tikz}
    \]
\end{definition}

\begin{definition}[Preparation Correctness Property]
    Suppose we have a preparation gadget associated with a QECC with distance $2t + 1$.
    The gadget satisfies the \emph{FT Preparation Correctness Property}, abbreviated \emph{PCP}, if, whenever $s \leq t$,
    \[
        \tikzfig{figures/Gottesman/PCP.tikz}
    \]
\end{definition}

\begin{definition}[FT Measurement Correctness Property]
    Suppose we have a preparation gadget associated with a QECC with distance $2t + 1$.
    The gadget satisfies the \emph{FT Measurement Correctness Property}, abbreviated \emph{MCP}, if, whenever $r + s \leq t$,
    \[
        \tikzfig{figures/Gottesman/MCP.tikz}
    \]
\end{definition}

\subsection{Fault-tolerant Quantum Error Correction}

\begin{definition}
    A fault-tolerant error correction gadget (which I will usually abbreviate FTEC in the future) is represented by the following picture:
    \[
        \tikzfig{figures/Gottesman/QEC-Box.tikz}
    \]
    The horizontal lines represent blocks of the QECC.
\end{definition}

\begin{definition}[FT Error Correction Recovery Property]
    Suppose we have an error correction gadget associated with a QECC with distance $2t + 1$.
    The gadget satisfies the \emph{FT Error Correction Recovery Property}, abbreviated \emph{ECRP}, if, whenever $s \leq t$,
    \[
        \tikzfig{figures/Gottesman/ECRP.tikz}
    \]
\end{definition}
\begin{definition}[FT Error Correction Correctness Property]
    Suppose we have an error correction gadget associated with a QECC with distance $2t + 1$.
    The gadget satisfies the \emph{FT Error Correction Correctness Property}, abbreviated \emph{ECCP}, if, whenever $r + s \leq t$,
    \[
        \tikzfig{figures/Gottesman/ECCP.tikz}
    \]
\end{definition}

\printbibliography

\end{document}